\documentclass[11pt]{article}
\usepackage[a4paper,margin=1.9cm]{geometry}
\usepackage[T1]{fontenc}
\usepackage{textcomp}
\usepackage[spanish]{babel}
\usepackage{amsmath,amssymb}
\usepackage{lmodern}
\renewcommand{\familydefault}{\sfdefault}
\usepackage{enumitem}

\newcommand{\vect}[2]{\begin{pmatrix}#1\\#2\end{pmatrix}}
\usepackage{tikz}
\usetikzlibrary{calc,arrows.meta,patterns,decorations.pathmorphing}
\usepackage[table]{xcolor}
\usepackage{multicol}
\usepackage{parskip}

\definecolor{unalblue}{RGB}{31,73,124}

\setlist[enumerate,1]{itemsep=1.4em,topsep=1em}

\newcommand{\examheader}{%
  \noindent
  \begin{minipage}[t]{0.6\linewidth}
    {\small\bfseries Geometría Vectorial y Analítica --- Parcial 1}\\
    {\small Semestre 2014--01}
  \end{minipage}%
  \hfill
  \begin{minipage}[t]{0.35\linewidth}
    \raggedleft\small Escuela de Matemáticas. Universidad Nacional de Colombia.
  \end{minipage}\\[2pt]
  \rule{\linewidth}{0.6pt}
}
\newcommand{\examfooter}{%
  \vfill
  \rule{\linewidth}{0.4pt}\\[1pt]
  {\scriptsize\textit{Transcripción digital --- MathUNAL}\hfill\textcolor{unalblue}{\textbf{\thepage}}}
}
\pagestyle{empty}

\newcommand{\nota}[1]{{\small\itshape\color{unalblue!70!black} #1}}

\begin{document}
\examheader

\begin{center}
{\Large Primer Examen Parcial de Geometría Vectorial}\\[4pt]
{\large TEMA B}\\[4pt]
{\small 22 de marzo de 2014}
\end{center}

\vspace{10pt}
\textbf{PUNTAJE. SOLO PARA USO OFICIAL}

\begin{center}
\begin{tabular}{|c|c|c|c|c|c|c|c|c|}
\hline
1 & 2(a) & 2(b) & 3(a) & 3(b) & 4 & 5 & 6 & TOTAL \\
\hline
 & & & & & & & & \\
\hline
\end{tabular}
\end{center}

\vspace{10pt}
\textbf{INSTRUCCIONES}

No está permitido el uso de calculadora ni de otros aparatos electrónicos. Solucione las preguntas en los espacios en blanco asignados a cada una de ellas. \textbf{Duración: 1 hora y 40 minutos.} No está permitido sacar hojas en blanco ni ningún tipo de apuntes durante el examen. Utilice la página 6 para realizar sus operaciones en borrador.

\begin{enumerate}
\item (12\%) Considere las rectas $\mathcal{L}_1: x+(2s)y-3=0$ y $\mathcal{L}_2: x+(3s)y=0$ con $s\in\mathbb{R}$. ¿Hay algún valor de $s$ para el cual las rectas $\mathcal{L}_1$ y $\mathcal{L}_2$ son perpendiculares?. ¿Hay algún valor de $s$ para el cual dichas rectas son paralelas? Justifique sus respuestas.

\item Sean $A=\vect{2}{-1}$, $B=\vect{-1}{1}$ y $C=\vect{-1}{-1}$.

\begin{enumerate}[label=(\alph*)]
\item (12\%) Halle la distancia del punto $C$ a la recta que pasa por los puntos $A$ y $B$.
\item (12\%) Calcule el área del triángulo con vértices $A$, $B$ y $C$.
\end{enumerate}

\item Sean $U=\vect{1}{2}$, $V=\vect{-3}{1}$ y $W=\vect{1}{-1}$.

\begin{enumerate}[label=(\alph*)]
\item (12\%) Muestre que $U$ y $V$ son linealmente independientes.
\item (12\%) Encuentre los escalares $r$ y $s$ tales que $W=rU+sV$.
\end{enumerate}

\item (14\%) Sean $U$ y $V$ vectores no nulos ortogonales. Pruebe que para todo vector $X$ del plano, $X = \mathrm{P}_U(X) + \mathrm{P}_V(X)$, donde $\mathrm{P}_U(X)$ y $\mathrm{P}_V(X)$ son, respectivamente, las proyecciones de $X$ sobre los vectores $U$ y $V$.

\item (12\%) Sea $\mathcal{L}$ la recta con ecuaciones paramétricas $x=3+t$, $y=5+t$, y sea $\mathcal{L}'$ la recta con ecuación vectorial $X=\vect{2}{-1}+t\vect{-1}{2}$, $t\in\mathbb{R}$. Encuentre el ángulo dirigido de la recta $\mathcal{L}$ a la recta $\mathcal{L}'$.

\item (14\%) Sean $U$ y $V$ vectores no nulos ortogonales, Pruebe que para todo vector $X$ del plano, $X = \mathrm{P}_U(X) + \mathrm{P}_V(X)$, donde $\mathrm{P}_U(X)$ y $\mathrm{P}_V(X)$ son, respectivamente, las proyecciones de $X$ sobre los vectores $U$ y $V$.

\end{enumerate}

\examfooter
\end{document}
